% g6e_K.tex — 6thGradeReviewExplanations, pp. 43–52:
% Section M (Decimal Computation) M1–M6 and Section N
% (Fraction Operations And Properties) N1–N7.
% Fragment: \input into a wrapper that loads mathdocs.sty.

% ===========================================================================
\lesson{M. Decimal Computation}

% ---------------------------------------------------------------------------
\question{M1}

\blockhead{Question}
\noindent $428.5 + 37.26 = $\ansblank\hint{decimal}

\checked{$465.76$}

\blockhead{What the question is asking}
\noindent This question asks you to add $428.5$ and $37.26$ (numbers that use a
dot for parts of a whole).

\blockhead{Important words}
\begin{words}
  \item A \term{decimal}uses a decimal point to show tenths and hundredths.
  \item To \term{add decimals}\unskip, line up the decimal points so equal place
        values are in the same column.
  \item The \term{sum}is the answer to an addition problem.
\end{words}

\blockhead{Work the problem one tiny step at a time}
\begin{steps}
  \item Write $428.5$ as $428.50$ to match hundredths.
  \item Stack the numbers with decimal points lined up (see the grid below).
  \item Add hundredths: $0 + 6 = 6$.
  \item Add tenths: $5 + 2 = 7$.
  \item Add ones: $8 + 7 = 15$. Write $5$ in ones and carry $1$ to the tens.
  \item Add tens: $2 + 3 + 1$ (carry) $= 6$.
  \item Add hundreds: $4 + 0 = 4$.
  \item Sum: $465.76$.
\end{steps}

\begin{center}
{\setlength{\tabcolsep}{0.15em}%
\begin{tabular}{cc}
      & $428.50$ \\
  $+$ & $37.26$  \\ \cline{1-2}
      & $465.76$ \\
\end{tabular}}
\end{center}

\finalans{$465.76$}

\blockhead{Why this makes sense}
\noindent Estimate: $429 + 37 = 466$, close to $465.76$.

% ---------------------------------------------------------------------------
\question{M2}

\blockhead{Question}
\noindent $700.3 - 284.67 = $\ansblank\hint{decimal}

\checked{$415.63$}

\blockhead{What the question is asking}
\noindent This question asks you to subtract $284.67$ from $700.3$ and write the
answer with a dot for parts of a whole.

\blockhead{Important words}
\begin{words}
  \item A \term{decimal}uses a decimal point to show tenths and hundredths.
  \item To \term{subtract decimals}\unskip, line up the decimal points so equal
        place values are in the same column.
  \item The \term{difference}is the answer to a subtraction problem.
  \item \term{Break apart}means split the number being subtracted into
        place-value parts and subtract one part at a time.
\end{words}

\blockhead{Work the problem one tiny step at a time}
\begin{steps}
  \item Write $700.3$ with two decimal places: $700.30$.
  \item Stack $700.30$ above $284.67$ with the decimal points lined up (see the
        grid below).
  \item Break $284.67$ into parts: $200 + 80 + 4 + 0.67$. Subtract one part at a
        time from $700.30$.
  \item Subtract the hundreds: $700.30 - 200 = 500.30$.
  \item Subtract the tens: $500.30 - 80 = 420.30$.
  \item Subtract the ones: $420.30 - 4 = 416.30$.
  \item Subtract six tenths: $416.30 - 0.60 = 415.70$.
  \item Subtract seven hundredths: $415.70 - 0.07 = 415.63$.
  \item So the difference is $415.63$.
\end{steps}

\begin{center}
{\setlength{\tabcolsep}{0.15em}%
\begin{tabular}{cc}
      & $700.30$ \\
  $-$ & $284.67$ \\ \cline{1-2}
      & $415.63$ \\
\end{tabular}}
\end{center}

\finalans{$415.63$}

\blockhead{Why this makes sense}
\noindent Check by adding the difference to what you subtracted:
$415.63 + 284.67$. Hundredths: $3 + 7 = 10$ (write $0$, carry $1$).
Tenths: $6 + 6 + 1 = 13$ (write $3$, carry $1$).
Ones: $5 + 4 + 1 = 10$ (write $0$, carry $1$).
Tens: $1 + 8 + 1 = 10$ (write $0$, carry $1$).
Hundreds: $4 + 2 + 1 = 7$. Result $700.30$, which matches the starting number.
Check passes.

% ---------------------------------------------------------------------------
\question{M3}

\blockhead{Question}
\noindent $18.4 \times 12 = $\ansblank\hint{decimal}

\checked{$220.8$}

\blockhead{What the question is asking}
\noindent This question asks you to multiply $18.4$ by $12$.

\blockhead{Important words}
\begin{words}
  \item To \term{multiply decimals}\unskip: first multiply as if the numbers were
        whole numbers.
  \item Then place the decimal point by counting all digits after the decimal
        points in the factors.
  \item A \term{decimal digit}is a digit to the right of the decimal point.
  \item A \term{factor}is a number being multiplied. Here the factors are $18.4$
        and $12$.
  \item The \term{product}is the answer to a multiplication problem.
\end{words}

\blockhead{Work the problem one tiny step at a time}
\begin{steps}
  \item Ignore the point: multiply $184 \times 12$.
  \item $184 \times 10 = 1840$.
  \item $184 \times 2 = 368$.
  \item Add: $1840 + 368 = 2208$.
  \item Count decimal digits: $18.4$ has $1$ decimal digit; $12$ has $0$ decimal
        digits; total $1$ decimal place.
  \item Place the point one digit from the right: $220.8$.
\end{steps}

\finalans{$220.8$}

\blockhead{Why this makes sense}
\noindent Estimate: $18 \times 12 = 216$, close to $220.8$.

% ---------------------------------------------------------------------------
\question{M4}

\blockhead{Question}
\noindent $168.8 \div 8 = $\ansblank\hint{decimal}

\checked{$21.1$}

\blockhead{What the question is asking}
\noindent This question asks you to divide $168.8$ by $8$.

\blockhead{Important words}
\begin{words}
  \item The \term{dividend}is the number being divided ($168.8$).
  \item The \term{divisor}is the number you divide by ($8$).
  \item The \term{quotient}is the answer to a division problem.
  \item When dividing a decimal by a whole number, divide as usual and keep the
        decimal point in the same place in the quotient (line it up under the
        dividend's decimal point).
  \item To \term{bring down}a digit means move the next digit into the current
        leftover so you can keep dividing.
\end{words}

\blockhead{Work the problem one tiny step at a time}
\begin{steps}
  \item Divide $8$ into $16$: $8 \times 2 = 16$. Write $2$ in the tens place of
        the quotient.
  \item Bring down $8$: $8 \div 8 = 1$. Write $1$ in the ones place.
  \item Place the decimal point in the quotient directly above the decimal point
        in $168.8$.
  \item Bring down $8$ (tenths): $8 \div 8 = 1$. Write $1$ in the tenths place.
  \item Quotient: $21.1$.
  \item Check step: $21.1 \times 8 = 168.8$ (because $21 \times 8 = 168$ and
        $0.1 \times 8 = 0.8$).
\end{steps}

\begin{center}
{\setlength{\tabcolsep}{0.18em}%
\begin{tabular}{ccccccc}
      &     &     & $2$ & $1$ & .   & $1$ \\ \cline{3-7}
  $8$ & $)$ & $1$ & $6$ & $8$ & .   & $8$ \\
\end{tabular}}
\end{center}

\begin{center}
Quotient digits sit above matching places; the decimal points line up.
\end{center}

\finalans{$21.1$}

\blockhead{Why this makes sense}
\noindent Check: $21.1 \times 8$. $21 \times 8 = 168$ and $0.1 \times 8 = 0.8$,
total $168.8$. Correct.

% ---------------------------------------------------------------------------
\question{M5}

\blockhead{Question}
\noindent $4.2 \times 1.5 = $\ansblank\hint{decimal}

\checked{$6.3$}

\blockhead{What the question is asking}
\noindent This question asks you to multiply $4.2$ by $1.5$.

\blockhead{Important words}
\begin{words}
  \item A \term{factor}is a number being multiplied. Here the factors are $4.2$
        and $1.5$.
  \item The \term{product}is the answer to a multiplication problem.
  \item To \term{multiply decimals}\unskip: multiply as if they were whole
        numbers, then place the decimal point by counting all decimal digits in
        the factors.
  \item A \term{decimal digit}is a digit to the right of the decimal point.
\end{words}

\blockhead{Work the problem one tiny step at a time}
\begin{steps}
  \item Ignore the points for a moment: multiply $42 \times 15$.
  \item $42 \times 10 = 420$.
  \item $42 \times 5 = 210$.
  \item Add: $420 + 210 = 630$.
  \item Count decimal digits: $4.2$ has $1$, $1.5$ has $1$, total $2$.
  \item Place the point two places from the right in $630$: $6.30$.
  \item Drop the trailing zero: $6.30 = 6.3$.
\end{steps}

\finalans{$6.3$}

\blockhead{Why this makes sense}
\noindent Estimate: $4 \times 1.5 = 6$, close to $6.3$.

% ---------------------------------------------------------------------------
\question{M6}

\blockhead{Question}
\noindent $7.2 \div 0.8 = $\ansblank

\checked{$9$}

\blockhead{What the question is asking}
\noindent This question asks you to divide $7.2$ by $0.8$.

\blockhead{Important words}
\begin{words}
  \item The \term{dividend}is the number being divided ($7.2$).
  \item The \term{divisor}is the number you divide by ($0.8$).
  \item The \term{quotient}is the answer to a division problem.
  \item A \term{power of ten}is a number like $10$, $100$, or $1000$ made by
        multiplying tens together.
  \item To \term{divide by a decimal}\unskip, multiply the dividend and the
        divisor by the same power of ten so the divisor becomes a whole number.
\end{words}

\blockhead{Work the problem one tiny step at a time}
\begin{steps}
  \item Name the parts: dividend $7.2$, divisor $0.8$.
  \item Multiply both by $10$ (one power of ten) to clear one decimal place in
        the divisor: $7.2 \times 10 = 72$ and $0.8 \times 10 = 8$.
  \item Now compute $72 \div 8$.
  \item $8 \times 9 = 72$, so $72 \div 8 = 9$.
  \item Therefore $7.2 \div 0.8 = 9$.
  \item Check: $0.8 \times 9 = 7.2$.
\end{steps}

\finalans{$9$}

\blockhead{Why this makes sense}
\noindent Check: $0.8 \times 9 = 7.2$. Yes.

% ===========================================================================
\lesson{N. Fraction Operations And Properties}

% ---------------------------------------------------------------------------
\question{N1}

\blockhead{Question}
\noindent $\dfrac{1}{2} + \dfrac{1}{4} = $\ansblank\hint{simplest fraction}

\checked{$\dfrac{3}{4}$}

\blockhead{What the question is asking}
\noindent This question asks you to add $\dfrac{1}{2}$ and $\dfrac{1}{4}$, then
shrink the answer if something bigger than $1$ divides both the top and the
bottom evenly.

\blockhead{Important words}
\begin{words}
  \item A \term{common denominator}is a shared bottom number so fractions can be
        added.
  \item To add fractions with the same denominator, add the numerators and keep
        the denominator.
\end{words}

\blockhead{Work the problem one tiny step at a time}
\begin{steps}
  \item The denominators are $2$ and $4$. A common denominator is $4$.
  \item Change $\dfrac{1}{2}$ into fourths: multiply top and bottom by $2$:
        $\dfrac{1 \times 2}{2 \times 2} = \dfrac{2}{4}$.
  \item Now add: $\dfrac{2}{4} + \dfrac{1}{4} = \dfrac{2 + 1}{4} = \dfrac{3}{4}$.
  \item $\dfrac{3}{4}$ is already in simplest terms.
\end{steps}

\finalans{$\dfrac{3}{4}$}

\blockhead{Why this makes sense}
\noindent One half is two fourths. Two fourths plus one fourth is three fourths.

% ---------------------------------------------------------------------------
\question{N2}

\blockhead{Question}
\noindent $\dfrac{5}{6} - \dfrac{1}{3} = $\ansblank\hint{simplest fraction}

\checked{$\dfrac{1}{2}$}

\blockhead{What the question is asking}
\noindent This question asks you to subtract $\dfrac{1}{3}$ from $\dfrac{5}{6}$,
then shrink the answer if needed.

\blockhead{Important words}
\begin{words}
  \item A \term{common denominator}is a shared bottom number so fractions can be
        subtracted.
  \item To \term{simplify}a fraction means divide top and bottom by the same
        common factor until they share no common factor greater than $1$.
  \item To subtract fractions with the same denominator, subtract the numerators
        and keep the denominator.
\end{words}

\blockhead{Work the problem one tiny step at a time}
\begin{steps}
  \item Denominators $6$ and $3$. Common denominator: $6$.
  \item Change $\dfrac{1}{3}$: multiply top and bottom by $2$: $\dfrac{2}{6}$.
  \item Subtract: $\dfrac{5}{6} - \dfrac{2}{6} = \dfrac{3}{6}$.
  \item Simplify $\dfrac{3}{6}$: write the division facts $3 \div 3 = 1$ and
        $6 \div 3 = 2$.
  \item So $\dfrac{3}{6} = \dfrac{1}{2}$.
\end{steps}

\finalans{$\dfrac{1}{2}$}

\blockhead{Why this makes sense}
\noindent Five sixths take away two sixths leaves three sixths, which is one
half.

% ---------------------------------------------------------------------------
\question{N3}

\blockhead{Question}
\noindent $\dfrac{5}{8} \times \dfrac{4}{15} = $\ansblank\hint{simplest fraction}

\checked{$\dfrac{1}{6}$}

\blockhead{What the question is asking}
\noindent This question asks you to multiply the equal-parts-over-a-whole
amounts and shrink the answer until nothing bigger than $1$ divides both the top
and the bottom evenly.

\blockhead{Important words}
\begin{words}
  \item To \term{multiply fractions}\unskip: multiply numerators; multiply
        denominators. Then simplify.
  \item To \term{cancel}means divide out the same factor from a numerator and a
        denominator before multiplying.
\end{words}

\blockhead{Work the problem one tiny step at a time}
\begin{steps}
  \item Write $\dfrac{5 \times 4}{8 \times 15} = \dfrac{20}{120}$.
  \item Simplify $\dfrac{20}{120}$: write $20 \div 20 = 1$ and
        $120 \div 20 = 6$, so $\dfrac{1}{6}$.
  \item Or cancel first: $4$ and $8$ share a factor of $4$, so divide those out:
        $\dfrac{5}{8} \times \dfrac{4}{15} = \dfrac{5}{2} \times \dfrac{1}{15}$.
  \item Multiply: $\dfrac{5 \times 1}{2 \times 15} = \dfrac{5}{30}$.
  \item Also cancel: $5$ and $30$ share a factor of $5$:
        $\dfrac{5}{30} = \dfrac{1}{6}$ (because $5 \div 5 = 1$ and
        $30 \div 5 = 6$).
\end{steps}

\finalans{$\dfrac{1}{6}$}

\blockhead{Why this makes sense}
\noindent Cross-check: $5 \times 4 = 20$ and $8 \times 15 = 120$, and
$\dfrac{20}{120} = \dfrac{1}{6}$.

% ---------------------------------------------------------------------------
\question{N4}

\blockhead{Question}
\noindent $\dfrac{7}{8} \div \dfrac{1}{4} = $\ansblank\hint{simplest fraction}

\checked{$\dfrac{7}{2}$}

\blockhead{What the question is asking}
\noindent This question asks you to divide $\dfrac{7}{8}$ by $\dfrac{1}{4}$,
then shrink the answer if needed.

\blockhead{Important words}
\begin{words}
  \item The \term{reciprocal}of a fraction flips the numerator and the
        denominator (swap top$\leftrightarrow$bottom). Example: the reciprocal
        of $\dfrac{1}{4}$ is $\dfrac{4}{1}$.
  \item To \term{flip}a fraction means the same swap: top$\leftrightarrow$bottom.
  \item To \term{divide by a fraction}\unskip, multiply by its reciprocal.
  \item To \term{cancel}means divide out the same factor from a numerator and a
        denominator before multiplying (same idea as N3).
\end{words}

\blockhead{Work the problem one tiny step at a time}
\begin{steps}
  \item Rewrite: $\dfrac{7}{8} \div \dfrac{1}{4} = \dfrac{7}{8} \times
        \dfrac{4}{1}$.
  \item Multiply: $\dfrac{7 \times 4}{8 \times 1} = \dfrac{28}{8}$.
  \item Simplify $\dfrac{28}{8}$: write $28 \div 4 = 7$ and $8 \div 4 = 2$, so
        $\dfrac{7}{2}$.
  \item Or cancel $4$ with $8$ first (same factor top and bottom):
        $\dfrac{7}{8} \times \dfrac{4}{1} = \dfrac{7}{2}$.
\end{steps}

\finalans{$\dfrac{7}{2}$}

\blockhead{Why this makes sense}
\noindent How many fourths fit into seven eighths? Multiplying by $4$ scales the
amount: $\dfrac{7}{2} = \mixed{3}{1}{2}$.

% ---------------------------------------------------------------------------
\question{N5}

\blockhead{Question}
\noindent Rewrite using the commutative property: $9 \times b = $

\checked{$b \times 9$}

\blockhead{What the question is asking}
\noindent This question asks you to rewrite $9 \times b$ by swapping the order of
the two numbers being multiplied (order can change; the answer stays the same).

\blockhead{Important words}
\begin{words}
  \item The \term{commutative property of multiplication}says
        $a \times b = b \times a$. You may switch the order of the factors; the
        product stays the same.
  \item A \term{factor}is a number or letter being multiplied.
\end{words}

\blockhead{Work the problem one tiny step at a time}
\begin{steps}
  \item Start with $9 \times b$.
  \item Name the factors in order: first factor $9$, second factor $b$.
  \item The commutative property lets you switch the order of the factors.
  \item After the switch, first factor is $b$, second factor is $9$: write
        $b \times 9$.
  \item Numeric check with a number for $b$: if $b = 3$, then $9 \times 3 = 27$
        and $3 \times 9 = 27$. Same product.
\end{steps}

\finalans{$b \times 9$}

\blockhead{Why this makes sense}
\noindent Order changes; the product meaning stays the same. This is not the
associative property (associative only moves parentheses).

% ---------------------------------------------------------------------------
\question{N6}

\blockhead{Question}
\noindent Rewrite using the associative property: $(x + 5) + 8 = $

\checked{$x + (5 + 8)$}

\blockhead{What the question is asking}
\noindent This question asks you to rewrite $(x + 5) + 8$ by moving the curved
marks to group the addition in a different way (grouping can change; the answer
stays the same).

\blockhead{Important words}
\begin{words}
  \item The \term{associative property of addition}says
        $(a + b) + c = a + (b + c)$. You may change the grouping; the sum stays
        the same.
  \item \term{Regroup}means move the parentheses without changing the
        left-to-right order of the terms.
\end{words}

\blockhead{Work the problem one tiny step at a time}
\begin{steps}
  \item Start with $(x + 5) + 8$. Here $x$ and $5$ are grouped first.
  \item List the three addends in order: $x$, then $5$, then $8$. Keep that
        order.
  \item Associative moves only the parentheses. Do not reorder into
        $8 + (x + 5)$.
  \item Move the parentheses to group $5$ and $8$: $x + (5 + 8)$.
  \item Numeric check with a number for $x$: if $x = 2$, then
        $(2 + 5) + 8 = 7 + 8 = 15$ and $2 + (5 + 8) = 2 + 13 = 15$. Same sum.
\end{steps}

\finalans{$x + (5 + 8)$}

\blockhead{Why this makes sense}
\noindent Associative regroups; it does not reorder. Writing $8 + (x + 5)$
equals the same sum, but the required associative rewrite keeps order and only
moves parentheses.

% ---------------------------------------------------------------------------
\question{N7}

\blockhead{Question}
\noindent Is subtraction associative? Compare $(12 - 4) - 3$ and $12 - (4 - 3)$.

\checked{No}

\blockhead{What the question is asking}
\noindent This question asks whether subtraction still gives the same answer when
you group the numbers in two different ways.

\blockhead{Important words}
\begin{words}
  \item \term{Associative}would mean regrouping never changes the answer.
  \item If the two results are different, subtraction is \textit{not}
        associative.
\end{words}

\blockhead{Work the problem one tiny step at a time}
\begin{steps}
  \item Compute the left grouping: $(12 - 4) - 3$.
  \item Inside first: $12 - 4 = 8$.
  \item Then $8 - 3 = 5$.
  \item Compute the right grouping: $12 - (4 - 3)$.
  \item Inside first: $4 - 3 = 1$.
  \item Then $12 - 1 = 11$.
  \item Compare: $5 \neq 11$.
  \item Because the answers differ, subtraction is not associative. Answer: No.
\end{steps}

\finalans{No}

\blockhead{Why this makes sense}
\noindent Addition is associative, but subtraction is not. Changing parentheses
in subtraction can change the result, as $5$ versus $11$ shows.
